\scanpage{118}
\uline{N\textsuperscript{o} 1}\quad \unsure{Préschémas} cpts et loc. cpts. Anneaux flasques.

\uline{N\textsuperscript{o} 2}\quad \deleted{Systèmes projectifs}\inserted{Proschémas} flasques \deleted{\unsure{proschémas} et \unsure{espaces schématiques} \illegible{2} \unsure{flasques}}

espaces schématiques compacts

\uline{N\textsuperscript{o} 3}\quad \deleted{Espaces \unsure{indschématiques} \unsure{Ind-compacts}} Espaces schématiques.

[Catégorie \unsure{recollable} pour \deleted{\unsure{tous}} les rec. ouverts...]

\margin{Morphismes \unsure{quasi-compacts} \unsure{projectifs} propres}

\uline{N\textsuperscript{o} 4}\quad Opérations sommes, \deleted{\unsure{lim}} et $\underleftarrow{\lim}$ \deleted{finies} des espaces schématiques

\inserted{\unsure{cas non quasi-compact $=$ réunion des compacts}}

\uline{N\textsuperscript{o} 5}\quad \deleted{\unsure{Ens quasi} \unsure{schématiques}} Espaces quasi-schématiques .....

\deleted{\diagram{\begin{tikzcd}[ampersand replacement=\&,column sep=large]
\left(\begin{gathered}\text{Préschémas quasi-compacts}\\\text{quasi-séparés}\end{gathered}\right)
 \arrow[r,"\alpha"]
\&\left(\begin{gathered}\text{Esp. schématiques}\\\text{compacts}\end{gathered}\right)
 \arrow[r,"{\substack{\text{\unsure{équiv.}}\\\deleted{\text{\unsure{naturelle}}}\\\simeq}}"]
\&\left(\operatorname{Pro}(\mathrm{Sch})\ \inserted{\unsure{\lim}}\ \deleted{\unsure{\mathrm{finis}}}\ \text{q. cpts}\right)
 \arrow[r,"{\substack{\unsure{\gamma}\\\text{fid,}\ \deleted{\text{\unsure{surjectif}}}\\\inserted{\text{\unsure{submersif}}}}}"]
\&\left(\begin{gathered}\text{schémas compacts}\\\text{réduits}\end{gathered}\right)
 \arrow[l,bend left=30,"{\sigma\quad\text{pl. fid.}}"']
 \arrow[lll,bend left=40,"{\deleted{\text{\illegible{1}}}\quad\text{pl. fid.}}"']
\end{tikzcd}}}

\uline{N\textsuperscript{o} 6}\quad Pb à \unsure{poser}. Espaces de Néron!

\diagram{\begin{tikzcd}[ampersand replacement=\&,column sep=huge]
\begin{gathered}(\text{Préschémas quasi-compacts}\\\text{quasi-séparés})\\X\end{gathered}
 \arrow[r,"\alpha","{\substack{\text{\unsure{submersif}}\\\text{[fidèle modulo}\\\text{\unsure{nilpotents}].}}}"']
 \arrow[rr,bend left=35,"\sigma"]
\&\begin{gathered}(\mathrm{Cons\ cpt})\\\text{''}\\(\operatorname{Pro}(\mathrm{Sch})\ \text{flasques})\\\mathcal Y\end{gathered}
 \arrow[r,"{\beta=\underleftarrow{\lim}}","{\substack{\text{\unsure{submersif}}\\\text{fidèle}}}"']
\&\begin{gathered}(\text{schémas compacts}\\\text{réduits})\\K\end{gathered}
 \arrow[l,bend left=35,"{\beta'\quad\text{pl. fidèle}}"']
 \arrow[ll,bend left=48,"{\sigma'\quad\text{pl. fid.}\quad\unsure{\text{imm. ouv.}}}"']
\end{tikzcd}}

\margin{$\unsure{\rho|\sigma}$ [pleinement fidèle quand on est ``de type fini'' sur un anneau fixé....]}

\margin{$\beta'(K)=$ système proj. trivial défini par $K$}

\[
\operatorname{Hom}(K,\sigma X)\xrightarrow{\sim}\operatorname{Hom}(\sigma'K,X)
\]
\[
\operatorname{Hom}(K,\beta\mathcal Y)\xrightarrow{\sim}\operatorname{Hom}(\beta'K,\mathcal Y)
\]
\[
\operatorname{Hom}(K,\beta(\alpha X))=\operatorname{Hom}(\beta'K,\alpha X)
\]
\margin{117}
